\documentclass{template}

\usepackage{mathtools}

\begin{document}

\section{Meeting \#1 Notes}

\subsubsection*{Exercises}

\begin{problem}
    Assume we are given the elliptic curve $E(\mathbb{K})$ in affine form and
    denote by $\overline{E}(\mathbb{K})$ the corresponding projective closure.
    What is $\#\overline{E}(\mathbb{K})\setminus E(\mathbb{K})$?
\end{problem}

\textbf{Solution.} We need to find such points on $\overline{E}(\mathbb{K})$
that are not included in $E(\mathbb{K})$. Recall that corresponding 
equations have the following form:
\begin{align*}
    E(\mathbb{K}): & \quad y^2 = x^3 + ax + b, \\
    \overline{E}(\mathbb{K}): & \quad Y^2Z = X^3 + aXZ^2 + bZ^3,
\end{align*}

where $a,b \in \mathbb{K}$. Notice that given $(X:Y:Z) \in
\overline{E}(\mathbb{K})$, we can almost always map it to $E(\mathbb{K})$ by
using projection map $\pi: (X:Y:Z) \mapsto (X/Z,Y/Z)$. The only issue we
encounter is when $Z=0$: so $E(\mathbb{K})$ does not contain any points on
$\overline{E}(\mathbb{K})$ with $Z=0$. What are these points? Substitute $Z=0$
into $\overline{E}(\mathbb{K})$: $X^3=0$, so $X=0$. Note, however, that $Y$ is
arbitrary, so up to equivalence, we have the so-called \textit{point at
infinity} $(0:1:0)$. Thus we conclude $\#\left(\overline{E}(\mathbb{K})\setminus
E(\mathbb{K})\right)=1$.

\begin{problem}
    Prove that $\mathbb{F}_q^{\times}$ is cyclic (where $q=p^n$ generally for
    prime $p$).
\end{problem}

\textbf{Proof.} For that, we prove the auxiliary lemma.

\begin{lemma}
    Let $H$ be a finite group of order $n$. Suppose that, for all divisors $d$
    of $n$, the set $\Omega_d=\{x \in H: x^d=1\}$ has at most $d$ elements. Then
    $H$ is cyclic.
\end{lemma}

\textit{Proof.} If there exists $x \in H$ of order $d$, the group $\langle x
\rangle$ generated by $x$ is cyclic of order $d$. In view of the hypothesis, all
elements $y \in H$ such that $y^d=1$ belong to $\langle x \rangle$. In
particular, all elements of $H$ of order $d$ are generators of $\langle x
\rangle$, and these are in number $\varphi(d)$. Hence, the number of elemets of
$H$ of order $d$ is $0$ or $\varphi(d)$. It cannot be $0$ since otherwise
$\sum_{d \mid n}\varphi(d)$ would show that the number of elements in $H$ is
less than $n$, contrary to the hypothesis. In particular, there exists an
element $x \in H$ of order $n$ and $H=\langle x \rangle$.

To prove our original theorem, simply use this lemma with
$H=\mathbb{F}_q^{\times}$ and $n=q-1$. Obviously, equation $x^d=1$ has at most
$d$ solutions over $\mathbb{F}_q$.

\section*{Quadratic Reciprocity Law}

The primary protocol where Quadratic Reciprocity Law is used 
is \textbf{Goldwasser-Micali cryptosystem}. 

$\mathsf{KeyGen}(1^{\lambda}) \to (\mathsf{pk}, \mathsf{sk})$. One generates two
large primes $p$ and $q$, and computes $N = pq$. Then, one finds some
non-residue $x$ such that $\left(\frac{x}{p}\right) = \left(\frac{x}{q}\right) =
-1$. Then, $\mathsf{pk} \gets (x, N)$, while $\mathsf{sk} \gets (p,q)$.

$\mathsf{Enc}(\mathsf{pk}, \mathbf{m} \in \{0,1\}^{\ell}) \to \mathbf{c} \in
\mathbb{Z}_N^{\ell}$. Upon receiving a binary string $\mathbf{m}$ of length
$\ell$, one generates the ciphertext $\mathbf{c}$ as follows: $c_i \gets
y_i^2x^{m_i}$ over $\mathbb{Z}_N$ where $y_i \leftarrow_R
\mathbb{Z}_N^{\times}$.

$\mathsf{Dec}(\mathsf{sk}, \mathbf{c} \in \mathbb{Z}_N^{\ell}) \to \mathbf{m} \in
\{0,1\}^{\ell}$. For each $i$, using the prime factorization $(p,q)$, one
determines whether the value $c_i$ is a quadratic residue. If so, $m_i=0$,
otherwise $m_i=1$.

\textbf{Why this is secure?} Given $x \in \mathbb{Z}_N$, there is no effective
way to determine whether $x$ is a quadratic non-residue, given no factorization 
of $N$.

\textbf{Why given $p,q$ we are good?} Because $x$ is a quadratic residue in
$\mathbb{Z}_N$ iff it is both quadratic residue in $\mathbb{Z}_p$ and
$\mathbb{Z}_q$. Thus, it suffices to check whether both $x^{(p-1)/2}=1$ and
$x^{(q-1)/2}=1$.
 
\section{Meeting \#3 Notes}

\begin{problem}
    Assume $f(\underline{X}) \in \mathbb{Z}[X_1,\dots,X_n]$ is a homogeneous 
    polynomial in $n$ variables of degree $n$. Let $\alpha$ be the 
    coefficient of $(X_1\dots X_n)^{p-1}$ in $f(\underline{X})^{p-1}$. Prove 
    that
    \begin{equation*}
        \#\{\underline{X} \in \mathbb{F}_p^n: f(\underline{X}) = 0\} = (-1)^{n+1}\alpha \; \text{in} \; \mathbb{F}_p
    \end{equation*}

    Consider the special case of elliptic curves.
\end{problem}

\textbf{Proof.} Denote by $V := \{\underline{X} \in \mathbb{F}_p^n:
f(\underline{X}) = 0\}$. Define $g(X) := 1 - X^{p-1}$. This function is
obviously the characteristic function of $\{0\}$. In such case, we can count
zero points simply by:
\begin{equation*}
    \# V = \sum_{\underline{X} \in \mathbb{F}_p^n} g \circ f(\underline{X}) = \sum_{\underline{X} \in \mathbb{F}_p^n} (1-f(\underline{X})^{p-1}) = -\sum_{\underline{X} \in \mathbb{F}_p^n}f(\underline{X})^{p-1}
\end{equation*}

Let $f(\underline{X})^{p-1} = \sum_{\underline{u}}c_{\underline{u}}\underline{X}^{\underline{u}}$. Then, we have:
\begin{align*}
    \# V &= -\sum_{\underline{X} \in \mathbb{F}_p^n}\sum_{\underline{u}}c_{\underline{u}}\underline{X}^{\underline{u}} = -\sum_{\underline{u}}\sum_{\underline{X} \in \mathbb{F}_p^n}c_{\underline{u}}\underline{X}^{\underline{u}} = -\sum_{\underline{u}}\sum_{\underline{X} \in \mathbb{F}_p^n}c_{\underline{u}}\prod_{j=1}^n X_j^{u_j} \\
    &= -\sum_{\underline{u}}c_{\underline{u}}\prod_{j=1}^n \sum_{X_j \in \mathbb{F}_p}X_j^{u_j} = -\sum_{\underline{u}}c_{\underline{u}}\prod_{j=1}^n S(X^{u_j}),
\end{align*}

where $S(\cdot)$ is the function we used during Chevalley–Warning theorem proof.
For the given $\underline{u}$, for product to be non-zero, every $u_j$ must be
divisible by $p-1$. Thus, the degree of the monomial must be at least $n(p-1)$.
However, note that $\deg f(\underline{X})^{p-1} \leq (p-1)n$, so the only term 
that contributes is in fact $u_1=u_2=\dots=u_n=p-1$. Therefore,
\begin{equation*}
    \# V = (-1)^{n+1}c_{p-1,\dots,p-1} = (-1)^{n+1}\alpha
\end{equation*}

\textbf{Application to Elliptic Curve.} For Elliptic Curve, $f(X,Y,Z) = Y^2Z -
X^3 - aXZ^2 - bZ^3$. In particular, $\#\{(X,Y,Z) \in \mathbb{F}_p^3:
f(X,Y,Z)=0\} = \alpha$.

\newpage

\section{Meeting \#4 Notes}

These notes are based on~\cite[Section 1]{ireland}. In this section, we define
the unique factorization domain (UFD) and Euclidean domain, and show why
$\mathbb{Z}$ and $k[x]$ are UFD and Euclidean domains.

\begin{definition}[Euclidean domain]
    We say that $R$ is an \textbf{Euclidean domain} if there is a function
    $\lambda: R \setminus \{0\} \to \mathbb{Z}_{\geq 0}$ such that if $a,b \in
    R$, $b \neq 0$, there exist $q,r \in R$ such that $a = qb + d$ and either
    $d=0$ or $\lambda(d) < \lambda(b)$.
\end{definition}

Further, we introduce the following notation: for some elements $a_1,\dots,a_n
\in R$ in ring $R$, let $\langle a_1,\dots,a_n \rangle = \{\sum_{i=1}^n r_ia_i:
r_i \in R\}$. We say $I$ is finitely generated if there are $a_1,\dots,a_n$ such
that $I = \langle a_1,\dots,a_n \rangle$. If $I = \langle a \rangle$ for some $a
\in I$, we say that $I$ is a principal ideal. We say that $R$ is a principal
ideal domain (PID) if every ideal $I \subseteq R$ is principal.

\begin{theorem}
    Euclidean domain is PID.
\end{theorem}

\textbf{Proof.} Suppose $R$ is an Euclidean domain and $I \subseteq R$ is an
ideal. We shall prove that $I = \langle a \rangle$ for some $a \in I$. Let us
construct it as follows: $a$ is such that $\lambda(a)$ is the smallest in the
set $\{\lambda(b): b \in I \setminus \{0\}\}$. We prove that such choice
suffices. First notice that $\langle a \rangle \subseteq I$ is due to the
definition of ideal. To prove $I \subseteq \langle a \rangle$, we take arbitrary
$b \in I$, $b \neq 0$. Divide $b$ by $a$: there exist $q,r \in R$ such that $b =
aq+r$ with $r=0$ or $\lambda(r) < \lambda(a)$. Note that $r=b-aq \in I$, thus
$\lambda(r) < \lambda(a)$ is not possible due to the choice of $a$. Thus, $r=0$
and thus $b=aq \in \langle a \rangle$, so $I \subseteq \langle a \rangle$.

Further discussion now related to PID only. We introduce several new notions.
\textit{(i)} $b$ divides $a$ iff exists $c \in R$ such that $a=bc$ or $\langle b
\rangle \subseteq \langle a \rangle$ \textit{(ii)} $u \in R$ is a unit iff $u
\mid 1$ or $\langle u \rangle = R$ \textit{(iii)} $a,b \in R$ are associates if
$a=bu$ for some unit $u \in R$ or $\langle a \rangle = \langle b \rangle$,
\textit{(iv)} $p \in R$ is irreducible if $a \mid p$ implies $a$ is a unit or an
associate of $p$, \textit{(v)} $p \in R$ is prime if $p \neq 0$ and $p \mid ab$
implies $p \mid a$ or $p \mid b$.

\begin{definition}
    $d$ is said to be \textbf{greatest common divisor (GCD)} of two 
    elements $a,b \in R$ if:
    \begin{itemize}
        \item[\textit{(a)}] $d \mid a$ and $d \mid b$.
        \item[\textit{(b)}] $d' \mid a$ and $d' \mid b$ implies $d' \mid d$.
    \end{itemize}
\end{definition}

\begin{proposition}[Generalization of Bezout Identity]
    Let $R$ be a PID and $a,b \in R$. Then $a$ and $b$ have a greatest 
    common divisor $d$ and $\langle a,b \rangle = \langle d \rangle$.
\end{proposition}

\textbf{Proof.} Let us consider the ideal $\langle a,b \rangle$. Since $R$ is a
PID, there exists some $d$ such that $\langle a,b \rangle = \langle d \rangle$.
We prove that $d$ is a GCD. First notice that $\langle a \rangle \subseteq
\langle a,b \rangle = \langle d \rangle$ and similarly $\langle b \rangle
\subseteq \langle d \rangle$, which translates to $d \mid a$ and $d \mid b$.
Thus, the first condition for GCD is satisfied. Now suppose $d' \mid a$ and $d'
\mid b$. Then, $\langle a \rangle, \langle b \rangle \subseteq \langle
d'\rangle$. Thus, $\langle d \rangle = \langle a, b \rangle \subseteq \langle d'
\rangle$, so $d' \mid d$. We are done.

\textbf{Corollary 1.} If $a,b \in R$ are relatively prime, $\langle a,b \rangle
= R$.

\textbf{Corollary 2.} If $R$ is a PID and $p \in R$ is irreducible, then $p$ is
prime.

\textbf{Proof.} Suppose $p \mid ab$ and $p$ does not divide $a$ (otherwise, we
are done). Then, $p$ and $a$ are relatively prime, so by Corollary 1 one has
$\langle a, p \rangle = R$. Thus, $\langle ab, pb \rangle = \langle b \rangle$.
Since $ab \in \langle p \rangle$ and $pb \in \langle p \rangle$, we have
$\langle b \rangle \subseteq \langle p \rangle$. Thus, $p \mid b$.

Now we show that every non-zero element of $R$ is a product of irreducible
elements. We do this in two steps: (i) if $a \in R$, $a \neq 0$, there is an
irreducible dividing $a$, (ii) $a$ is product of irreducibles.

\begin{lemma}\label{lemma:stabilizes}
    Let $\langle a_1 \rangle \subseteq \langle a_2 \rangle \subseteq \langle a_3
    \rangle \subseteq \dots$ be an ascending chain of ideals in PID $R$. Then
    the chain stabilizes in finitely many steps.
\end{lemma}

\textbf{Proof.} Let $I = \bigcup_{n \in \mathbb{N}} \langle a_n \rangle$. $I$ is
an ideal. Indeed, take any $b \in I$, then $b \in \langle a_k \rangle$ for some
$k$. Then, trivially $\langle b \rangle \subseteq \langle a_k \rangle \subseteq
I$. Since $R$ is a PID, $I = \langle a \rangle$ for some $a \in R$. For the same
argument as before, we get $I = \langle a \rangle \subseteq \langle a_k \rangle
\subseteq I$ for some $a_k$. Thus, $I = \langle a_k \rangle$, and the result
follows.

\begin{proposition}
    Every non-zero element of $R$ is a product of irreducibles.
\end{proposition}

\textbf{Proof.} Let $a \in R$, $a \neq 0$, $a$ is not a unit. We first show that
$a$ is divisible by some irreducible element. If $a$ is irreducible, we are
done. Otherwise, $a=a_1b_1$ where $a_1,b_1$ are non-units. If $a_1$ is
irreducible, we are done. Otherwise, $a_1=a_2b_2$ where $a_2,b_2$ are non-units.
If $a_2$ is irreducible, we are done. Otherwise continue as before. Notice that
$\langle a \rangle \subseteq \langle a_1 \rangle \subseteq \langle a_2 \rangle
\subseteq \dots$. By Lemma~\ref{lemma:stabilizes}, this stabilizes after some
$a_k$, which is irreducible.

We now show that $a$ is the product of irreducibles. If $a$ is irreducible, we
are done. Otherwise let $p_1 \mid a$ with $p_1$ --- irreducible. Then
$a=p_1c_1$. If $c_1$ is a unit, we are done. Otherwise, let $p_2$ be an
irreducible such that $p_2 \mid c_1$. Then, $a=p_1p_2c_2$. If $c_2$ is a unit,
we are done. Otherwise, continue as before. Notice that $\langle a \rangle
\subseteq \langle c_1 \rangle \subseteq \langle c_2 \rangle \subseteq \dots$. By
Lemma~\ref{lemma:stabilizes}, this stabilizes, so from some $c_k$, we have
units. Since $p_kc_k$ is irreducible, we are done.

Now we define the order function.

\begin{definition}
    Define the order of $a$ relatively to prime $p$ (denoted as $\text{ord}_pa$)
    as the largest integer $n$ such that $p^n \mid a$, but not $p^{n+1} \mid a$.
\end{definition}

\begin{lemma}
    Order always exists.
\end{lemma}

\textbf{Proof.} Construct the converse. That is, for each $m>0$, one has
$a=b_mp^m$. Then, $pb_{m+1}=b_m$ so that $\langle b_1 \rangle \subseteq \langle
b_2 \rangle \subseteq \langle b_3 \rangle \subseteq \dots$ is an infinite
ascending chain of ideals, which is impossible due to
Lemma~\ref{lemma:stabilizes}.

\begin{lemma}
    If $a,b \in R$ with $a,b \neq 0$, then $\text{ord}_p(ab) = \text{ord}_pa +
    \text{ord}_pb$. 
\end{lemma}

\textbf{Proof.} Let $\alpha = \text{ord}_pa$ and $\beta=\text{ord}_pb$. Then,
$a=p^{\alpha}c$ and $b=p^{\beta}d$ with $p$ not dividing $c,d$. Note that 
$ab=p^{\alpha+\beta}cd$. Since $p$ is prime, $p$ does not divide $cd$. 

Let $S$ be a set of primes in $R$ with the following two properties:
\begin{itemize}
    \item[\textit{(a)}] Every prime in $R$ is associate to a prime in $S$.
    \item[\textit{(b)}] No two primes in $S$ are associate. 
\end{itemize}

\textbf{Example.} In $\mathbb{Z}$, we choose $S$ to be the set of positive
primes. In $S$, we choose $S$ to be the set of monic irreducible polynomials.

\begin{theorem}
    Let $R$ be a PID and $S$ a set of primes with the properties given above.
    Then if $a \in R$, $a \neq 0$, we can write
    \begin{equation*}
        a = u\prod_{p \in S}p^{e(p)},
    \end{equation*}
    where $u$ is a unit. The unit $u$ and the exponents $e(p)$ are uniquely
    determined by $a$. In fact, $e(p) = \text{ord}_pa$.
\end{theorem}
\textbf{Proof.} Existence follows immediately from aforementioned discussion.

Take $\text{ord}_q(\cdot)$ to both sides with $q \in S$. Then,
\begin{equation*}
    \text{ord}_qa = \text{ord}_qu + \sum_{p \in S}e(p)\text{ord}_qp
\end{equation*}

Now, $\text{ord}_q(u)=0$ from the definition, and $\text{ord}_qp = 1$ iff $q=p$.
Thus, $\text{ord}_qa = e(q)$. Since the exponents $e(q)$ are uniquely determined
so is the unit $u$. This completes the proof.

\begin{table}
\centering
\begin{tabular}{cccc}
\hline
\textbf{Property} & $\mathbb{Z}$ & $K[x]$ & \textbf{General Euclidean $R$} \\
\hline
\textbf{Unique factorization} & $a=\pm\prod p^{e(p)}$ & $f(x)=c\prod p(x)^{e(p)}$ & $a=u\prod_{p \in S}p^{e(p)}$ \\
\hline
\textbf{Euclidean function $\lambda$} & $|\cdot|$ & $\deg(\cdot)$ & --- \\
\hline
\textbf{Ideals} & $n\mathbb{Z}$ & $\langle f(x) \rangle$ & Every ideal is principal \\
\hline
\textbf{Divisibility} & $a \mid b \iff b=ac$ & $f \mid g \iff g=fh$ & $a \mid b \iff \langle b \rangle \subseteq \langle a \rangle$ \\
\hline
\textbf{Units} & $\{\pm 1\}$ & $K^{\times}$ & $u$ s.t. $\langle u \rangle = R$ \\
\hline
\textbf{Associates} & $a = \pm b$ & $f \sim g \iff f=cg$ & $a \sim b \iff \langle a \rangle = \langle b \rangle$ \\
\hline
\textbf{Irreducible elements} & $2,3,5,7,\dots$ & $x^2+1$ over $\mathbb{R}$ & $a \mid p \implies \langle a \rangle = \langle p \rangle$ or $R$ \\
\hline
\textbf{GCD (Bezout)} & $\gcd(a,b)=ua+vb$ & $\gcd(f,g)=uf+vg$ & $d$ s.t. $\langle a,b \rangle = \langle d \rangle$ \\
\hline
\textbf{Canonical primes $S$} & Positive primes & Monic irreducibles & Representative of associate class \\
\hline
\end{tabular}
\caption{Comparison between $\mathbb{Z}$, $k[x]$, and a general integral domain/PID.}
\end{table}

\subsection{Other Euclidean Domains}

\begin{theorem}
    Let $\mathbb{Z}[i] = \{a+bi: a,b \in \mathbb{Z}\}$ be the ring of Gaussian
    integers. Then, it is an Euclidean domain for $\lambda(z) = z\overline{z}$.
\end{theorem}

\textbf{Proof.} For $\alpha = a+bi$ and $\gamma=c+di$, we have $\alpha/\gamma =
r+si$ for some rational $r,s$. Choose integers $m,n \in \mathbb{Z}$ such that
$|r-m|\leq \frac{1}{2}$ and $|s-n| \leq \frac{1}{2}$ (simply by rounding). Set
$\delta := m+ni$. Clearly, $\delta \in \mathbb{Z}[i]$. Moreover,
$\lambda(\frac{\alpha}{\gamma} - \delta)=(r-m)^2+(s-n)^2 \leq \frac{1}{2}$. Set
$\rho := \alpha - \gamma\delta$. Then, $\rho \in \mathbb{Z}[i]$ and either
$\rho=0$ or $\lambda(\rho) =
\lambda(\gamma(\frac{\alpha}{\gamma}-\delta))=\lambda(\gamma)\lambda(\frac{\alpha}{\gamma}-\delta)
\leq \frac{1}{2}\lambda(\gamma)<\lambda(\gamma)$.

\begin{theorem}
    Let $\mathbb{Z}[\omega] = \{a+b\omega: \omega^2+\omega+1=0\}$ be the ring of
    Eisenstein integers. Then, $\mathbb{Z}[\omega]$ is an Euclidean domain with
    $\lambda(z) = z\overline{z}$.
\end{theorem}

\newpage
\section{Meeting \#6 Notes}

\subsection{Algebraic Numbers and Integers}

\begin{definition}
    An \textbf{algebraic number} is a complex number $\alpha$ which is a root of
    some polynomial with rational coefficients $f \in \mathbb{Q}[x], f \neq 0$.
    \textbf{Algebraic integer} is a complex number $\omega$ which is a root of
    some polynomial with integer coefficients $g \in \mathbb{Z}[x], g \neq 0$.
\end{definition}

\begin{proposition}
    A rational number $r \in \mathbb{Q}$ is an algebraic integer iff $r \in
    \mathbb{Z}$.
\end{proposition}

\textbf{Proof.} Clearly integer $r \in \mathbb{Z}$ is an algebraic integer
(since it is a root of $x-r$). Now suppose $r \in \mathbb{Q}$ is an algebraic
integer, thus it is a root of some $\sum_{j=0}^n a_jx^j$, thus $\sum_{j=0}^n
a_jr^j = 0$ with $a_j \in \mathbb{Z}$. Since $r$ is rational, assume
$r=\frac{p}{q}$. Thus $q^n\sum_{j=0}^n a_j\frac{p^j}{q^j} = 0$ or $\sum_{j=0}^n
a_jq^{n-j}p^j = 0$. Note that $q$ divides all terms but last, so $q \mid p^j$.
Since $\text{gcd}(p,q)=1$, this implies $q \mid p$, thus $r \in \mathbb{Z}$.

\begin{definition}
    A subset $V \subseteq \mathbb{C}$ of the complex numbers is called a
    $\mathbb{Q}$-module if it is a finite dimensional space over $\mathbb{Q}$.
    Denote by $[\gamma_1,\gamma_2,\dots,\gamma_{\ell}]$ the set of all
    expressions $\sum_{i=1}^{\ell} r_i\gamma_i$ with $r_i \in \mathbb{Q}$ and
    $\gamma_i \in \mathbb{C}$.
\end{definition}

\begin{proposition}
    Let $V=[\gamma_1,\dots,\gamma_{\ell}]$ and suppose that $\alpha \in
    \mathbb{C}$ has the property that $\alpha\gamma \in V$ for all $\gamma \in
    V$. Then, $\alpha$ is an algebraic number.
\end{proposition}

\textbf{Proof.} Since $\alpha\gamma_i \in V$, we have $\alpha\gamma_i =
\sum_{j=1}^{\ell}a_{i,j}\gamma_j$ or simply $\alpha\boldsymbol{\gamma} =
A\boldsymbol{\gamma}$, which implies $\det(A - \alpha E_{\ell \times \ell}) =
0$. Since $A \in \mathbb{Q}^{\ell \times \ell}$, $\alpha$ is a root of some
polynomial with rational coefficients, and thus an algebraic number.

\begin{proposition}
    The set of algebraic numbers forms a field.
\end{proposition}

\textbf{Proof.} Suppose $\alpha_1$ and $\alpha_2$ are algebraic numbers. We
shall show that $\alpha_1+\alpha_2$ and $\alpha_1\alpha_2$ are algebraic numbers
as well. Suppose $\alpha_1^n + r_{n-1}\alpha_1^{n-1} + \dots + r_0 = 0$ and
$\alpha_2^m + s_{m-1}\alpha_2^{m-1} + \dots + s_0 = 0$, where all $r_i,s_i \in
\mathbb{Q}$. Let $V := [\alpha_1^i\alpha_2^j]_{(i,j) \in [n] \times [m]}$ be the
$\mathbb{Q}$-module. For $\gamma \in V$ let us show that $\alpha_1\gamma \in V$
and $\alpha_2\gamma \in V$. Indeed, suppose $\gamma =
\sum_{i,j}c_{i,j}\alpha_1^i\alpha_2^j$, then $\alpha_1\gamma =
\sum_{i,j}c_{i,j}\alpha_1^{i+1}\alpha_2^j$. It would be clearly in $V$ if we did
not have $\alpha_1^n$ term, which is not included in $V$ by definition. However,
$\alpha_1^n \in V$ since $\alpha_1^n = -\sum_{j \in [n]} r_j\alpha_1^j$, thus
$c_{n,j}\alpha_1^n\alpha_2^j \in V$, concluding the argument. 

Therefore, $(\alpha_1+\alpha_2)\gamma \in V$ and $(\alpha_1\alpha_2)\gamma \in
V$. By proposition above, this implies that both $\alpha_1+\alpha_2$ and 
$\alpha_1\alpha_2$ are algebraic numbers.

Finally, if $\alpha_1$ is an algebraic number, $\alpha_1^{-1}$ is algebraic
number as well. Indeed, $1+r_{n-1}\alpha_1^{-1} + \dots + r_0\alpha_1^{-n} = 0$. 

Note that proof for algebraic integers is the same, but we have a ring
structure.

Let $\Omega$ be the ring of algebraic integers. If $\omega_1,\omega_2,\gamma \in
\Omega$, we say that $\omega_1 \equiv \omega_2 \pmod{\gamma}$ if
$\omega_1-\omega_2 = \gamma\alpha$ for some $\alpha \in \Omega$. 

\begin{proposition}
    If $\omega_1,\omega_2 \in \Omega$ and $p \in \mathbb{Z}$ is prime, then
    \begin{equation*}
        (\omega_1 + \omega_2)^p \equiv \omega_1^p + \omega_2^p \pmod{p}
    \end{equation*}
\end{proposition}

\textbf{Proof.} $(\omega_1+\omega_2)^p = \sum_{j=0}^p \begin{pmatrix} p \\ j
\end{pmatrix} \omega_1^j\omega_2^{p-j}$. Again, $p$ divides all terms but
the first and last, so we reach the conclusion.

\begin{proposition}
    If $\alpha$ is an algebraic number then $\alpha$ is the root of a unique 
    monic irreducible $f \in \mathbb{Q}[x]$. Furthermore, if $g \in \mathbb{Q}[x]$,
    $g(\alpha)=0$, then $f \mid g$.
\end{proposition}

\textbf{Proof.} Let $f$ be any monic irreducible with $f(\alpha)=0$ (it exists
because $\alpha$ is the root of some $q \in \mathbb{Q}[x]$, and if $q$ is not
irreducible, we can decompose it into factors where $\alpha$ will also be a
root). Let us prove the second assertion first. If $f \mid g$ does not hold,
then $\text{gcd}(f,g) = 1$, and so $f(x)h(x) + g(x)t(x) = 1$ for some $h,t \in
\mathbb{Q}[x]$. Setting $x=\alpha$ gives the contradiction. Uniqueness follows 
immediately.

The polynomial defined in proposition is called the \textbf{minimal polynomial}
of $\alpha$. $\alpha$ is called the algebraic number of degree $n$ if 
corresponding $f$ has a degree $n$. Now, define the field:
\begin{equation*}
    \mathbb{Q}(\alpha) = \{g(\alpha)/h(\alpha): g, h \in \mathbb{Q}[x], h(\alpha) \neq 0\}
\end{equation*}

\begin{proposition}
    If $\alpha \in \Omega$, then $\mathbb{Q}(\alpha) \cong \mathbb{Q}[\alpha]$.
\end{proposition}

\textbf{Proof.} Clearly $\mathbb{Q}[\alpha] \subseteq \mathbb{Q}(\alpha)$. If
$h(\alpha) \in \mathbb{Q}[\alpha]$, $h(\alpha) \neq 0$, then $f(x)$ does not
divide $h(x)$ where $f(x)$ is the minimal polynomial of $\alpha$. Then,
$\text{gcd}(h(x),f(x)) = 1$ and so $s(x)f(x)+t(x)h(x)=1$ for some $s,t \in
\mathbb{Q}[x]$. Set $\alpha=1$, thus $t(\alpha)h(\alpha)=1$, so $h(\alpha) \in
\mathbb{Q}[\alpha]^{\times}$. If $\beta \in \mathbb{Q}(\alpha)$, then $\beta =
g(\alpha)h(\alpha)^{-1}$ and thus $\beta \in \mathbb{Q}[\alpha]$ since
$h(\alpha)$ is a unit.

\textbf{Corollary.} If $\alpha$ is of degree $n$, then
$[\mathbb{Q}(\alpha):\mathbb{Q}]=n$.

\subsection{The Quadratic Character of 2}

We shall prove the following fact differently.

\begin{proposition}
    $(2/p)=(-1)^{\varepsilon}$ where $\varepsilon \equiv (p^2-1)/8 \pmod{2}$.
\end{proposition}

\textbf{Proof.} Let $\zeta = e^{2\pi i/8}$ and let $\tau = \zeta+\zeta^{-1}$.
Clearly, both $\zeta$ and $\tau$ are algebraic integers, so we might work with
congruences in the ring of algebraic integers. Let $p$ be the odd prime in
$\mathbb{Z}$ and notice that
\begin{equation*}
    \tau^{p-1} = (\tau^2)^{(p-1)/2} = 2^{(p-1)/2} \equiv (2/p) \pmod{p}
\end{equation*}

It thus follows that $\tau^p \equiv (2/p)\tau \pmod{p}$. Here, we used the fact
that $\tau^2 = (\zeta+\zeta^{-1})^2=2$, and the last congruence is due to
Euler's criterion. By earlier proposition, $\tau^p \equiv \zeta^p + \zeta^{-p}
\pmod{p}$. But since $\zeta^8=1$, we have $\zeta^p+\zeta^{-p}=\zeta+\zeta^{-1}$
for $p \equiv \pm 1 \pmod{8}$ and $\zeta^p+\zeta^{-p} = \zeta^3 + \zeta^{-3} = -(\zeta+\zeta^{-1})$ 
for $p \equiv \pm 3 \pmod{8}$. Summarizing,
\begin{equation*}
    \zeta^p + \zeta^{-p} = \begin{cases}
        \tau, & \text{if} \; p \equiv \pm 1 \pmod{8}, \\
        -\tau, & \text{if} \; p \equiv \pm 3 \pmod{8}.
    \end{cases}
\end{equation*}

Thus, substituting into $\tau^p \equiv (2/p)\tau \pmod{p}$ we get
$(-1)^{\varepsilon}\tau \equiv (2/p)\tau \pmod{p}$ with $\varepsilon \equiv
\frac{p^2-8}{1} \pmod{3}$. Multiplying both sides by $\tau$ we get
$(-1)^{\varepsilon}2 \equiv (2/p)2 \pmod{p}$, but dividing both sides by $2$ we
finally get $(2/p)=(-1)^{\varepsilon}$, which is a desired result.

\subsection{Quadratic Gauss Sums}

We shall generalize the relation $(\zeta+\zeta^{-1})^2=2$ for arbitrary odd
prime $p$. Further, denote by $\zeta=e^{2\pi i/p}$ a primitive $p$th root of
unity (root of $x^p+1=0$).

\begin{lemma}
    $\sum_{t=0}^{p-1}\zeta^{at}$ equals $p$ if $a \equiv 0 \pmod{p}$. Otherwise, it is $0$.
\end{lemma}

\textbf{Proof.} If $a \equiv 0 \pmod{p}$, then $\zeta^a=1$, so the result
follows. If $a \not\equiv 0 \pmod{p}$, then $\zeta^a \neq 1$, and so we get the
geometric series sum $\sum_{t=0}^{p-1}\zeta^{at} =
\frac{\zeta^{ap}-1}{\zeta^a-1}=0$.

\textbf{Corollary.} Denote $\delta(x,y)$ the function which is $1$ iff $x \equiv
y \pmod{p}$. Then,
\begin{equation*}
    \delta(x,y) = p^{-1}\sum_{t=0}^{p-1}\zeta^{t(x-y)}
\end{equation*}

\begin{lemma}
    $\sum_t (t/p)=0$ where $(t/p)$ is the Legendre symbol.
\end{lemma}

\textbf{Proof.} By definition, $(0/p)=0$. Of the remaining $p-1$ terms, the half
are $+1$, the half is $-1$, so the result follows.

\begin{definition}
    $g_a = \sum_t (t/p)\zeta^{at}$ is called the \textit{quadratic Gauss sum}.
\end{definition}

\begin{proposition}
    $g_a = (a/p)g_1$.
\end{proposition}

\textbf{Proof.} If $a \equiv 0 \pmod{p}$, then $\zeta^{at}=1$, and so $g_a =
\sum_t (t/p) = 0$ by Lemma above. Suppose in contrary that $a \not\equiv 0
\pmod{p}$. Then,
\begin{equation*}
    \left(\frac{a}{p}\right)g_a = \sum_t \left(\frac{at}{p}\right)\zeta^{at} = \sum_x \left(\frac{x}{p}\right)\zeta^x = g_1.
\end{equation*}

Here we used the fact that $x \mapsto ax$ is a bijection. But since $(a/p)^2=1$,
the result follows by multiplying both sides by $(a/p)$.

Further, denote $g_1$ by $g$. If follows from proposition above that 
$g^2=g_a^2$ for all $a \not\equiv 0 \pmod{p}$. 

\begin{proposition}
    $g^2 = (-1)^{(p-1)/2}p$.
\end{proposition}

\textbf{Proof.} The idea of the proof is to evaluate $\sum_a g_ag_{-a}$ in two ways.

If $a \not\equiv 0 \pmod{p}$, then $g_ag_{-a} = (a/p)(-a/p)g^2 = (-1/p)g^2$. Thus,
\begin{equation*}
    \sum_{a}g_ag_{-a} = \left(\frac{-1}{p}\right)(p-1)g^2
\end{equation*}

On the other hand,
\begin{equation*}
    g_ag_{-a} = \sum_x \sum_y \left(\frac{x}{p}\right)\left(\frac{y}{p}\right)\zeta^{a(x-y)}
\end{equation*}

Apply $\sum_a$ to both sides:
\begin{equation*}
    \sum_a g_ag_{-a} = \sum_x \sum_y \left(\frac{x}{p}\right)\left(\frac{y}{p}\right)\sum_a \zeta^{a(x-y)} = \sum_x \sum_y \left(\frac{x}{p}\right)\left(\frac{y}{p}\right) \delta(x,y)p = (p-1)p
\end{equation*}

Thus we arrive at the result.

Let $p^* = (-1)^{(p-1)/2}p$. The equation $g^2=p^*$ is the desired analog of the
equation $\tau^2=2$. Let $q\neq p$. We shall prove the following theorem.
\begin{theorem}
    $(p/q)(q/p) = (-1)^{(p-1)(q-1)/4}$.
\end{theorem}

\textbf{Proof.} Let $p^*=(-1)^{(p-1)/2}$, as above. Now note:
\begin{equation*}
    g^{q-1} = (g^2)^{(q-1)/2} = (p^*)^{(q-1)/2} \equiv (p^*/q) \pmod{q}
\end{equation*}

Thus, $g^q \equiv (p^*/q)g \pmod{q}$. Now note:
\begin{equation*}
    g^q = \left(\sum_p (t/p)\zeta^t\right)^q = \sum_p (t/p)^q \zeta^{qt} \equiv g_q \pmod{q}
\end{equation*}

Thus it follows that $g^q \equiv g_q \equiv (q/p)g \pmod{q}$ and so
\begin{equation*}
    (q/p)g = (p^*/q)g \pmod{q}
\end{equation*}

Multiply both sides by $g$ and use $g^2=p^*$:
\begin{equation*}
    (q/p)p^* = (p^*/q)p^* \pmod{q},
\end{equation*}

implying that $(q/p)=(p^*/q) \pmod{q}$ or simply $(q/p)=(p^*/q)$. This proves the 
theorem. To see this, note that:
\begin{equation*}
    (p^*/q) = (-1/q)^{(p-1)/2}(p/q) = (-1)^{(q-1)(p-1)/4}(p/q)
\end{equation*}

\section{Meeting \#8 Notes}
\subsection{Zeta Function}

\begin{definition}
    Let $V$ be a non-singular, projective variety over $K=\mathbb{F}_q$. Let
    $N_s$ be the set of points rational over $K_s = \mathbb{F}_{q^s}$. The
    \textbf{(Hasse-Weil) zeta function} of $V$ is the series given by:
    \begin{equation*}
        Z(V, T) = \exp\left(\sum_{s>1} \frac{N_sT^s}{s}\right)
    \end{equation*}
\end{definition}

In particular, \cite{ireland} considers $f \in \mathbb{F}_q[T_0,\dots,T_n]$ ---
a homogeneous polynomial, and considers projective hypersurface $V :=
\#\overline{H}_f(\mathbb{F}_{q^s}) \subseteq \mathbb{P}^n(\mathbb{F}_{q^s})$.

\subsection{Examples}

\begin{example}[Single Point Variety]
    Let $V$ be a one-element set. Then, $Z(T)=1/(1-T)$.
\end{example}

\begin{example}
    Let $V=\mathbb{A}_K^d$, then $N_s = q^{sd}$, and thus $Z_V(T) =
    1/(1-q^dT)$.
\end{example}

\begin{example}[Hyperplane]
    Consider the hyperplane at infinity:
    \begin{equation*}
        \overline{H}_f(\mathbb{F}_{q^s}) := \{(x_0,\dots,x_n) \in \mathbb{P}^{n}(\mathbb{F}_{q^s}): x_0 = 0\}.
    \end{equation*}
    The number of points is $N_s = q^{s(n-1)}+q^{s(n-2)}+\dots+q^s+1$. It follows that:
    \begin{equation*}
        \sum_{s=1}^{\infty} \frac{N_su^s}{s} = \sum_{m=0}^{n-1}\sum_{s=1}^{\infty}\frac{(q^mu)^s}{s} = -\sum_{m=0}^{n-1}\log(1-q^mu),
    \end{equation*}
    where we used the identity $\sum_{s=1}^{\infty}w^s/s = -\log(1-w)$. Thus,
    \begin{equation*}
        Z_f(u) = \prod_{m=0}^{n-1} \frac{1}{1-q^mu}
    \end{equation*}
\end{example}

\begin{example}
    Consider the hypersurface $-y_0^2+y_1^2+y_2^2+y_3^2=0$. It can be computed
    that the number of points on such surface is
    $N_s=q^{2s}+q^s+1+\chi_s(-1)q^s$. Thus:
    \begin{equation*}
        \sum_{s=1}^{\infty}\frac{N_su^s}{s} = \sum_{s=1}^{\infty}\frac{(q^su)^s}{s} + 2\sum_{s=1}^{\infty}\frac{(qu)^s}{s} + \sum_{s=1}^{\infty}\frac{u^s}{s}
    \end{equation*}

    If $-1$ is a square in $\mathbb{F}_q$, then $\chi_s(-1)=1$ for all $s$,
    otherwise $\chi_s(-1) = -1$ for $s$ odd and $\chi_s(-1)=1$ for all $s$ even.
    In the first case, $Z(u) = (1-q^2u)^{-1}(1-qu)^{-2}(1-u)^{-1}$, while in the
    second $Z(u) = (1-q^2u)^{-1}(1-qu)^{-1}(1+qu)^{-1}(1-u)^{-1}$.
\end{example}

\begin{example}
    For elliptic curve $E$ over $K=\mathbb{F}_q$, there are algebraic integers
    $\alpha,\beta \in \mathbb{C}$ such that:
    \begin{equation*}
        Z(E/K, T) = \frac{(1-\alpha T)(1-\beta T)}{(1-T)(1-qT)}
    \end{equation*}
    with $|\alpha| = |\beta| = \sqrt{q}$.
\end{example}

\begin{lemma}[Hasse's Bound]
    For any elliptic curve $E(\mathbb{F}_q)$, one has:
    \begin{equation*}
        |\#E(\mathbb{F}_q) - 1 + q| \leq 2\sqrt{q}
    \end{equation*}
\end{lemma}

\textbf{Proof.} In general, note:
\begin{equation*}
    \frac{TZ'(T)}{Z(T)} = \sum_{s > 0}N_sT^s,
\end{equation*}

so $Z(V,T)$ actually encodes the number of points on $V$. Now, for 
$E(\mathbb{F}_q)$, we know $Z(V,T)$ explicitly, so let us find its logarithm:
\begin{equation*}
    \log Z(V,T) = \log(1-\alpha T) + \log(1-\beta T) - \log(1-T) - \log(1-qT)
\end{equation*}

Find formal derivative of both sides:
\begin{equation*}
    \frac{Z'(T)}{Z(T)} = \frac{1}{1-T} + \frac{q}{1-qT} - \frac{\alpha}{1-\alpha T} - \frac{\beta}{1-\beta T}
\end{equation*}

Expand $\frac{1}{1-T}=\sum_{n \geq 0}T^n$, $\frac{q}{1-qT} = \sum_{n \geq 0}q^{n+1}T^n, \dots$ and 
use the fact that $Z'(T)/Z(T) = \sum_{s > 0}N_{s}T^{s-1}$.
We obtain
\begin{equation*}
    \#E(\mathbb{F}_{q^k}) = q^k + 1 - (\alpha^k + \beta^k)
\end{equation*}

Using that $|\alpha^k + \beta^k| \leq 2q^{k/2}$, we obtain:
\begin{equation*}
    |\#E(\mathbb{F}_{q^k}) - (q^k+1)| \leq 2\sqrt{q^k}
\end{equation*}

\subsection{Historical Notes}

Recall the fact that $\sum_{n>0}\frac{1}{n^s} = \prod_{p}\frac{1}{1-p^{-s}}$.

\begin{definition}
    Let $A$ be an algebra of finite type over $\mathbb{Z}$. This means that $A
    \cong \mathbb{Z}[T_1,\dots,T_m]/I$ for $I \subseteq
    \mathbb{Z}[T_1,\dots,T_m]$. Define:
    \begin{equation*}
        \zeta_A(s) = \prod_{\mathfrak{m} \subseteq A} \frac{1}{1-\#(A/\mathfrak{m})^{-s}}
    \end{equation*}
\end{definition}

\begin{example}
    For $A=\mathbb{Z}$, we have a standard Riehmann function.

    By taking $X := \text{Spec}(A)$, we get that $\mathfrak{m}=\{x\}$ are closed
    points and $A/\mathfrak{m}$ is a finite field of some order $q_x$. Then,
    \begin{equation*}
        \zeta_X(s) = \prod_{x \in \text{Spec}(A)}\frac{1}{1-q_x^{-s}}
    \end{equation*}
\end{example}

\subsection{Conjectures}
The Weil (F.K. Schmidt 1931) was able to prove that for $d:=\deg f$:
\begin{equation*}
    Z_f(T) = \frac{P(T)}{(1-T)(1-qT)}, \quad P \in \mathbb{Z}[T], \, \deg P = (d-1)(d-2)
\end{equation*}

If $C$ is a curve, $\deg P = 2g, g=\text{genus}(C)$.


\begin{proposition}[11.1.1~\cite{ireland}]
    The zeta function is rational iff there exist complex numbers $\alpha_i$ and $\beta_j$
    such that $N_s = \sum_j \beta_j^s - \sum_i \alpha_i^s$.
\end{proposition}

\textbf{Proof.} On one hand $TZ'(T)/Z(T) = \sum_{s > 0}N_sT^s$, on the other
substitute $Z(T) = \frac{\prod_i (1-\alpha_iT)}{\prod_j (1-\beta_jT)}$.

\textbf{Conjectures.} Let $V$ be a non-singular, projective variety of 
dimension $d$ defined over the $K=\mathbb{F}_q$.
\begin{enumerate}
    \item \textit{Rationality.} Zeta function is a rational function 
    with coefficients in $\mathbb{Z}$ in $T := q^{-s}$:
    \begin{equation*}
        Z(V/K, T) = \frac{\prod_{j=1}^d P_{2j-1}(T)}{\prod_{j=1}^d P_{2j}(T)}, \quad P_j \in \mathbb{Z}[T], \; P_0 = 1-T, \; P_{2d} = 1-q^dT.
    \end{equation*}
    \item \textit{Riemann hypothesis.} The polynomials $P_1,\dots,P_{2d-1}$
    factor over $\mathbb{C}$ as $P_k(T) = \prod_j (1-\alpha_{k,j}T)$ with
    $|\alpha_{k,j}| = \sqrt{q^k}$.
    \item \textit{Functional equation.} Zeta function should satisfy:
    \begin{equation*}
        Z(V, n-s) = \pm q^{nE/2 - Es}Z(V, s),
    \end{equation*}
    where $E$ is the Euler characteristic of $V$.
    \item \textit{Betti numbers.} Degrees of $P_i$ are determined by cohomological 
    characterization of $V$.
\end{enumerate}

\subsection{Hasse-Davenport Relation}

\begin{theorem}
    Let $K := \mathbb{F}_q$, $K_s := \mathbb{F}_{q^s}$. Let $\chi' = \chi \circ
    N_{K_s/K}$ be the multiplicative character of $K_s$, while $\psi' = \psi
    \circ \text{Tr}_{K_s/K}$ be the additive character of $K_s$ (where typically
    $\psi'(\alpha) = \zeta_p^{\text{Tr}(\alpha)}$). Denote Gauss sum as
    $g(\chi,\psi) := \sum_{\alpha \in K} \chi(\alpha)\psi(\alpha)$. Then,
    \begin{equation*}
        (-1)^sg(\chi,\psi)^s = -g(\chi',\psi')
    \end{equation*}
\end{theorem}

We need some auxiliary lemmas to prove this theorem. Further for monic
polynomial $f(T) := T^n + \sum_{j=1}^n (-1)^jc_jT^{n-j}$, define $\lambda(f) :=
\psi(c_1)\chi(c_n)$.

\begin{lemma}
    $\lambda(fg) = \lambda(f)\lambda(g)$ for all monic $f,g \in k[T]$.
\end{lemma}

\textbf{Proof.} Let $f(T) = T^n + \sum_{i=1}^n (-1)^ia_iT^{n-i}$ and 
$g(T) = T^m + \sum_{j=1}^m (-1)^jb_jT^{m-j}$. Note:
\begin{equation*}
    f(T)g(T) = T^{n+m} + \sum_{k=1}^{n+m} (-1)^k c_k T^{n+m-k}.
\end{equation*}

We are interested only in $c_1$ and $c_{n+m}$. Note $c_1 = a_1 + b_1$ and
$c_{n+m} = b_ma_n$. Thus:
\begin{equation*}
    \lambda(fg) = \psi(a_1+b_1)\chi(b_ma_n) = \psi(a_1)\psi(b_1)\chi(b_m)\chi(a_n) = \lambda(f)\lambda(g)
\end{equation*}

\begin{lemma}
    Let $\alpha \in K_s$ and $f(T)$ be the monic irreducible polynomial for
    $\alpha$ over $K$. Then,
    \begin{equation*}
        \lambda(f)^{s/d} = \chi'(\alpha)\psi'(\alpha), \quad d = \deg f
    \end{equation*}
\end{lemma}

\textbf{Proof.} Note two things: if $K \subseteq K(\alpha) \subseteq K_s$ with
$[K(\alpha):K]=d$, then (a) $\text{Tr}_{K_s/K}(\alpha) = (s/d)c_1$ and (b)
$N_{K_s/K}(\alpha) = c_d^{s/d}$. Let us show them.

Recall $f(T) = (T-\alpha)\dots (T-\alpha^{q^{d-1}})$. Immediately
$\text{Tr}_{K(\alpha)/K} = c_1$ and $N_{K(\alpha)/K} = c_d$.
Finally 
\begin{equation*}
    \text{Tr}_{K_s/K}(\alpha) = \text{Tr}_{K(\alpha)/K}(\text{Tr}_{K_s/K(\alpha)}(\alpha)) = \text{Tr}_{K(\alpha)/K}([K_s:K(\alpha)]c_1) = \frac{s}{d}c_1
\end{equation*}
Similar logic can be applied to show $N_{K_s/K}(\alpha) = c_d^{s/d}$.

Now, recall that $\lambda(f) = \psi(c_1)\chi(c_d)$, thus
\begin{equation*}
    \lambda(f)^{s/d} = \psi(c_1)^{s/d}\chi(c_d)^{s/d} = \psi(c_1s/d)\chi(c_d^{s/d}) = \psi(\text{Tr}_{K_s/K}(\alpha))\chi(N_{K_s/K}(\alpha)) = \psi'(\alpha)\chi'(\alpha).
\end{equation*}

\begin{lemma}
    $g(\chi', \psi') = \sum_{f \in \mathcal{P}_s} (\deg f)\lambda(f)^{s/\deg f}$ where $f$ ranges
    over $\mathcal{P}_s$  --- all monic irreducible polynomials $f$ in $K[T]$ with $\deg f \mid s$.
\end{lemma}

\textbf{Proof.} By definition, $g(\chi', \psi') = \sum_{\alpha \in K_s} \chi'(\alpha)\psi'(\alpha)$. We partition the elements of $K_s$ based on their minimal polynomials over $K$.

Let $\mathcal{P}_s$ be the set of all monic irreducible polynomials in $K[T]$ whose degree divides $s$. For any $\alpha \in K_s$, let $f_\alpha \in K[T]$ be the minimal polynomial of $\alpha$. Since $K_s = \mathbb{F}_{q^s}$, we know that $\deg(f_\alpha)$ must divide $s$. Conversely, if $f \in \mathcal{P}_s$ with $d = \deg f$, then $f$ splits completely in $K_s$ and has exactly $d$ distinct roots in $K_s$.

Thus, we can rewrite the summation by grouping roots of the same polynomial:
\begin{equation*}
    \sum_{\alpha \in K_s} \chi'(\alpha)\psi'(\alpha) = \sum_{f \in \mathcal{P}_s} \sum_{\alpha: f(\alpha)=0} \chi'(\alpha)\psi'(\alpha).
\end{equation*}

Let $f \in \mathcal{P}_s$ with $d = \deg f$. By \textbf{Lemma 2}, for any root $\alpha$ of $f$, the value of the character product is constant:
\begin{equation*}
    \chi'(\alpha)\psi'(\alpha) = \lambda(f)^{s/d}.
\end{equation*}
Since $f$ has exactly $d$ roots in $K_s$, the inner sum evaluates to:
\begin{equation*}
    \sum_{\alpha: f(\alpha)=0} \chi'(\alpha)\psi'(\alpha) = d \cdot \lambda(f)^{s/d}.
\end{equation*}
Summing over all $f \in \mathcal{P}_s$, we obtain:
\begin{equation*}
    g(\chi', \psi') = \sum_{f \in \mathcal{P}_s} (\deg f)\lambda(f)^{s/\deg f}.
\end{equation*}

\textbf{Proof of Hasse-Davenport Relation.}
\begin{proposition}
    Let $\mathcal{M}(K)$ be all monic polynomials in $K[T]$. Then,
    \[\sum_{f \in \mathcal{M}(K)} \lambda(f)T^{\deg f} = \prod_{f \in \mathcal{P}(K)}(1-\lambda(f)T^{\deg f})^{-1}\]
\end{proposition}

For this identity not to look too random, it is analogous to the Euler product
expansion $\sum_{n \geq 1} a(n)n^{-s} = \prod_{p}(1-a(p)p^{-s})^{-1}$ for some
multiplicative bounded function $a$.

\textbf{Proof.} Expand into geometric series:
\begin{align*}
    \text{RHS} &= \prod_{f \in \mathcal{P}(K)}(1-\lambda(f)T^{\deg f})^{-1} \\
    &=\prod_{f \in \mathcal{P}(K)}\sum_{n \geq 0} \lambda(f)^{n}T^{n\deg f} \\
    &=\prod_{f \in \mathcal{P}(K)}\sum_{n \geq 0} \lambda(f^n)T^{\deg f^n} \\
    &=\sum_{f \in \mathcal{M}(K)} \lambda(f)T^{\deg f}
\end{align*}

The last equality more or less holds because each monomial in the former product is:
\begin{equation*}
    \prod_{f} \lambda(f^{n_f})T^{\sum_fn_f\deg f} = \lambda\left(\prod_f f^{n_f} \right)T^{\sum_f n_f\deg f}
\end{equation*}

Finally note that since $K[x]$ is a UFD, $\prod_f f^{n_f}$ corresponds to some
monic polynomial from $\mathcal{M}(K)$ of degree $\sum_f n_f\deg f$, thus we get
the desired result.

Now, for our $\lambda$, note that:
\begin{equation*}
    \sum_{f \in \mathcal{M}(K)} \lambda(f)T^{\deg f} = \sum_{n \geq 0}\sum_{\deg f=n}\lambda(f)T^{\deg f} 
\end{equation*}

For $n=0$, our sum is $1$. For $n=1$ we have:
\begin{equation*}
    \sum_{\deg f=1}\lambda(f) = \sum_{\alpha \in K}\lambda(T-\alpha) = \sum_{\alpha \in K}\chi(\alpha)\psi(\alpha) = g(\chi,\psi)
\end{equation*}

For $n>1$, we have:
\begin{equation*}
    \sum_{\deg f = n}\lambda(f) = \sum_{c_i \in K}\lambda(x^n - c_1x^{n-1} + \dots + (-1)^nc_n) = q^{n-2}\sum_{c_1,c_s \in K} \chi(c_s)\psi(c_1) = 0
\end{equation*}

Now take logarithmic derivative of both sides and multiply by $T$:
\begin{equation*}
    \frac{g(\chi,\psi)T}{1+g(\chi,\psi)T} = \sum_{f \in \mathcal{P}(K)}\frac{\lambda(f)(\deg f)T^{\deg f}}{1-\lambda(f)T^{\deg f}}
\end{equation*}

Expand geometric series:
\begin{equation*}
    \sum_{n \geq 1}(-1)^{n-1}g(\chi,\psi)^nT^n = \sum_{f \in \mathcal{P}(K)}\sum_{n \geq 1} (\deg f)\lambda(f)^nT^{n\deg f}
\end{equation*}

Equating coefficient of $T^s$:
\begin{equation*}
    (-1)^{s-1}g(\chi,\psi)^s = \sum_{\deg f \mid s} (\deg f)\lambda(f)^{s/\deg f}
\end{equation*}

By Lemma above, RHS is $g(\chi',\psi')$.

\subsection*{References}
\renewcommand{\refname}{\vspace{-1em}}
\bibliographystyle{alphaurl}
\bibliography{refs}

\end{document}